\chapter{CONCLUSÃO}
\label{cap:conclusao}

Toda sua experiência no LEMAF foi excepcional e de grande valia pro meu desenvolvimento profissional e pessoal, aprendendo a lidar com diversos problemas e pessoas.

A UFLA tive grande influencia na sua iniciativa no estagio, uma vez que a vaga me foi oferecida em um evento da SETI(Semana de Tecnologia da Informação) oferecido pela CompJunior\footnote{Empresa Junior do departamento de ciencia da computação} da UFLA e consegui a vaga após faler na entrevista sobre alguns projetos que já havia desenvolvido pra disciplinas,
 sendo o que chamou a atenção do entrevistar o projeto que foi desenvolvido graças a matéria de Desenvolvimento Mobile, um aplicativo de chat com integração com Firebase\footnote{Um Baas (Backend as a Service) para aplicações Web e Mobile do Google}.
Quando iniciei o estágio, meu conhecimento era somente relacionado ao contexto acadêmico lecionado pela universidade até o quinto período de graduação, o que em sua grande maioria eram somente conhecimentos e teorias sobre otimizações e códigos de baixo nível.
Após alguns meses como estagiário, já possuía propriedade intelectual e autoridade para opinar em decisões estruturais de projetos, tecnologias e até mesmo estabelecer prazos, obviamente após inúmeros erros e muito estudo.
 
O estagio foi complementar a sua formação na UFLA assim como a UFLA complementava o estagio, varíos problemas eram solucionados graças a materias lecionadas na mesma e varias materias que cursei se tornaram simples pois já havia experiência com aquilo.

Posso citar como exemplo as disciplinas de Banco de Dados, Sistemas Gerenciadores de Banco de Dados, Desenvolvimento Web, Modelagem e Implementação de Software, POO e muitas outras.

Os profissionais com os quais tive o prazer de trabalhar sempre foram muito agradáveis e solícitos, agregando muito valor ao meu desenvolvimento.

Devido aos problemas superados e projetos desenvolvidos no LEMAF, fui capaz de aprimorar minhas habilidades como desenvolvedor web e criar um bom portfólio, que me abriu diversas portas para o mercado de trabalho.
O estágio também me proporcionou uma visão abrangente de cargos e caminhos que poderei trilhar dentro da sua área de estudo, o que fez aumentar as oportunidades de escolhas e afirmar o que realmente tinha como desejo e vocação: ser desenvolvedor.

Posso afirmar com segurança que, os principais motivos da sua contratação na empresa Equals, meu emprego atual, foram minhas experiências com as tecnologias ReactJS, Spring e principalmente o desenvolvimento do Cria Design.

\begin{figure}[H]
\centering
\caption{Bilhete recebido da confraternização no fim de 2018} %legenda

\includegraphics[scale=0.3]{agradecimento}\\  % o 0.9 indica 90% do tamanho original
% pdfLaTeX aceita figuras no formato PNG, JPG ou PDF
% figuras vetoriais podem ser exportadas para eps e depois convertidas para pdf usando epstopdf
\label{fig:exemplo} %rotulo para refencia
\end{figure}

Toda sua trajetória no laboratório foi muito prazerosa e produtiva. No meu ultimo dia de estágio, fui convidado a participar do podcast da empresa, onde compartilhei um pouco da sua experiência com os outros funcionários \footnote{\url{http://blog.ti.lemaf.ufla.br/2019/08/22/podcast-5/}}.
